%=============================================================================
% TOEPLITZ TRUNCATION MATHEMATICS IN DFD's ALPHA DERIVATION
% Complete mathematical anatomy of how d=64 Toeplitz truncation works
% Extracted from v108 Unified Theory: section_alpha_locked.tex, appendix_K.tex,
% appendix_O_alpha57.tex, appendix_P_kalpha_MR.tex, appendix_X_neutrino.tex,
% and Alpha_Inverse_Formula.tex
%
% Created: 2026-03-22
%=============================================================================
\documentclass[12pt,a4paper]{article}
\usepackage{amsmath,amssymb,amsthm}
\usepackage{booktabs}
\usepackage{geometry}
\geometry{margin=2.5cm}
\usepackage{hyperref}
\usepackage{tcolorbox}

\newtheorem{theorem}{Theorem}
\newtheorem{lemma}{Lemma}
\newtheorem{corollary}{Corollary}
\newtheorem{definition}{Definition}

\title{The Toeplitz Truncation at $d=64$:\\[6pt]
\large Complete Mathematical Anatomy of DFD's $\alpha$ Derivation}
\author{Cross-Reference Document\\
Extracted from G.\ Alcock, \textit{Density Field Dynamics: A Complete Unified Theory} (v108)}
\date{March 2026}

\begin{document}
\maketitle

\begin{abstract}
This document extracts and organizes the complete mathematical structure of the
Toeplitz truncation at $d=64$ in the DFD derivation of $\alpha^{-1} = 137.03599985$.
We trace four specific mechanisms: (a)~how the Dirac/Laplacian spectrum is modified
by finite truncation, (b)~the traceless projection factor $(d-1)/d = 63/64$,
(c)~the regular-module boost $\varepsilon_{\mathrm{adj}} = d^2/(d^2-1) = 4096/4095$,
and (d)~the explicit connection to $\kappa_{U(1)}$ via the master formula
$\alpha^{-1} = 6\pi\,\kappa_{U(1)}$.
All source references are to the v108 Unified Theory.
\end{abstract}

\tableofcontents
\newpage

%=============================================================================
\section{Overview: What the Toeplitz Truncation Does}
\label{sec:overview}
%=============================================================================

The DFD internal geometry is $X = \mathbb{C}P^2 \times S^3$ (7-dimensional).
The continuous geometry has infinite-dimensional function spaces. The Toeplitz
truncation replaces the continuous $\mathbb{C}P^1 \subset \mathbb{C}P^2$ with
a \textbf{finite matrix algebra} $M_d(\mathbb{C})$, where $d = 64$.

This truncation has three independent mathematical consequences that enter the
$\alpha$ derivation:

\begin{enumerate}
\item \textbf{Spectral cutoff:} The finite-dimensional Hilbert space truncates
  the spectrum of the connection Laplacian, modifying the $a_4$ Seeley--DeWitt
  coefficient that yields the gauge kinetic term.

\item \textbf{Traceless projection} $(63/64)$: The determinant (trace) channel
  is removed from the $\mathrm{SU}(d)$ gauge field, projecting onto the
  traceless part and producing a factor $(d-1)/d$.

\item \textbf{Regular-module boost} $(4096/4095)$: The choice of finite Hilbert
  space as the algebra itself ($M_d(\mathbb{C})$, the ``regular module'') rather
  than a fundamental representation ($\mathbb{C}^d$) forces a trace-conversion
  factor $d^2/(d^2-1)$.
\end{enumerate}

%=============================================================================
\section{Origin of $d = 64$: The Spin$^c$ Structure on $\mathbb{C}P^2$}
\label{sec:origin-d}
%=============================================================================

The integer $d$ is not chosen---it is \textbf{derived} from the algebraic
geometry of $\mathbb{C}P^2$.

\subsection{The Canonical Bundle and Spin$^c$ Shift}

The canonical bundle of $\mathbb{C}P^2$ is:
\begin{equation}
K_{\mathbb{C}P^2} = \mathcal{O}(-3).
\end{equation}
The Spin$^c$ structure requires the determinant line bundle:
\begin{equation}
L_{\det} = K^{-1} = \mathcal{O}(3).
\end{equation}

\textbf{Source:} Section~8C (Sec.~\ref*{sec:alpha-finite-k} of the Unified Theory),
Table~XXVIII rows 1--2.

\subsection{Toeplitz Quantization Level}

The Toeplitz quantization is performed at level $m = k + 3$ on
$\mathbb{C}P^1 \subset \mathbb{C}P^2$. The shift by $+3$ arises because
restricting the line bundle $\mathcal{O}(k) \otimes L_{\det}$ to
$\mathbb{C}P^1$ yields $\mathcal{O}(k+3)$.

The dimension of the truncated Hilbert space is:
\begin{equation}
\boxed{d = \dim H^0(\mathbb{C}P^1, \mathcal{O}(m)) = m + 1 = k + 4.}
\end{equation}

With $k = k_{\max} = 60$ (from the Bridge Lemma / Spin$^c$ index):
\begin{equation}
d = 60 + 4 = 64.
\end{equation}

\textbf{Source:} Section~8C, Eq.~(8C.4) of the Unified Theory.

\subsection{Why $k_{\max} = 60$}

The maximum Chern--Simons level is a topological invariant:
\begin{equation}
k_{\max} = \chi(\mathbb{C}P^2, E) = \chi(\mathcal{O}(9)) + 5\chi(\mathcal{O})
= \binom{11}{2} + 5 = 55 + 5 = 60,
\end{equation}
where $E = \mathcal{O}(9) \oplus \mathcal{O}^{\oplus 5}$ is the twist bundle
fixed by minimal-padding of the SM hypercharge structure.

\textbf{Source:} Appendix~K, Eq.~(K.8).

%=============================================================================
\section{(a) How the Dirac/Laplacian Spectrum Is Modified}
\label{sec:spectrum}
%=============================================================================

\subsection{The Operator}

On the internal space $X = \mathbb{C}P^2 \times S^3$, DFD uses the
\textbf{connection Laplacian}:
\begin{equation}
P = -g^{ij}\nabla_i\nabla_j,
\end{equation}
acting on the internal bundle carrying emergent gauge degrees of freedom.
The U(1) factor is implemented via twisting by a line bundle over
$\mathbb{C}P^2$ with curvature proportional to the K\"ahler form $\omega$.

\textbf{Key property:} The K\"ahler form is parallel ($\nabla\omega = 0$),
so derivative terms in higher Seeley--DeWitt coefficients vanish automatically.
This is why this operator choice is ``locked''---alternative operators would
introduce ppm-level ambiguities.

\textbf{Source:} Section~8C, Sec.~\ref*{sec:alpha-operator-choice}.

\subsection{The Spectral Action}

The gauge kinetic term is extracted from the spectral action:
\begin{equation}
S = \mathrm{Tr}\, f(P/\Lambda^2),
\end{equation}
with a \textbf{plateau cutoff} function $f$ satisfying $f^{(n)}(0) = 0$ for
all $n \geq 1$. This forces all negative moments to vanish:
\begin{equation}
f_{-2} = f_{-4} = \cdots = 0,
\end{equation}
eliminating $a_6$ and higher heat-kernel contributions as potential tuning knobs.

\textbf{Source:} Section~8C, Sec.~\ref*{sec:alpha-regularization}.

\subsection{The Finite-$k$ Spectral Cutoff}

The Toeplitz truncation modifies the spectrum by imposing a
\textbf{finite-dimensional cutoff}. Instead of summing over the full infinite
spectrum of $P$ on $\mathbb{C}P^2$, the spectral action is restricted to the
$d$-dimensional truncated Hilbert space.

The determinant-channel removal at finite $d$ fixes the spectral cutoff as:
\begin{equation}
\boxed{\Lambda^3 = k \cdot \frac{d-1}{d} = k \cdot \frac{k+3}{k+4}
= 60 \times \frac{63}{64} = 59.0625.}
\end{equation}

This is the \emph{unique} finite-size factor permitted by the truncation rule.
It encodes the fact that in $d$ dimensions, the traceless gauge field has
$(d^2-1)$ generators rather than $d^2$, removing one degree of freedom
(the trace/determinant channel).

\textbf{Source:} Section~8C, Eq.~(8C.5); Alpha\_Inverse\_Formula.tex, Sec.~3.

\subsection{How This Enters the $a_4$ Coefficient}

The $a_4$ Seeley--DeWitt coefficient of the connection Laplacian on the
Toeplitz-truncated geometry determines the gauge kinetic term:
\begin{equation}
S_{\mathrm{gauge}} \supset -\frac{\kappa_{U(1)}}{2}\,\mathrm{Tr}\,
F^{(U(1))}_{ij} F^{(U(1))}_{ij},
\end{equation}
where $\kappa_{U(1)}$ is the U(1) stiffness coefficient computed from the $a_4$
trace over the internal fermion Hilbert space, incorporating the spectral cutoff
$\Lambda^3 = 59.0625$, the SM hypercharge weighting $w = 7/80$, the spectral
grading $g_F = 8$, and the microsector trace conversion factor.

\textbf{Source:} Alpha\_Inverse\_Formula.tex, Sec.~2--3.

%=============================================================================
\section{(b) Traceless Projection: $(d-1)/d = 63/64$}
\label{sec:traceless}
%=============================================================================

\subsection{Mathematical Origin}

The Toeplitz-truncated algebra is $M_d(\mathbb{C})$, which has $d^2$ real
generators. The gauge group $\mathrm{U}(d)$ decomposes as:
\begin{equation}
\mathrm{U}(d) = \mathrm{U}(1) \times \mathrm{SU}(d) / \mathbb{Z}_d.
\end{equation}

The determinant (trace) channel corresponds to the overall $\mathrm{U}(1)$ factor.
In the DFD construction, this channel must be \textbf{removed} because the
physical gauge fields are traceless (they live in $\mathfrak{su}(d)$, not
$\mathfrak{u}(d)$).

\subsection{Effect on the Spectral Cutoff}

The removal of the determinant channel reduces the effective number of modes
by one out of $d$:
\begin{equation}
\boxed{\text{Traceless projection factor} = \frac{d-1}{d} = \frac{63}{64}
= 0.984375.}
\end{equation}

This factor enters the spectral cutoff $\Lambda^3 = k(d-1)/d$ and thus
modifies the overall normalization of the $a_4$ coefficient. Physically, it
says that of the $d$ levels in the Toeplitz-truncated Hilbert space, only
$d-1 = 63$ contribute to the traceless gauge sector.

\subsection{Where It Appears in Table XXVIII}

In the derivation chain (Table~XXVIII of Section~8C):
\begin{center}
\begin{tabular}{lll}
\toprule
Component & Value & Source \\
\midrule
$(d-1)/d$ & $63/64$ & Traceless projection \\
\bottomrule
\end{tabular}
\end{center}

\textbf{Status:} Derived (not fitted). It is a consequence of the
$\mathrm{SU}(d)$ vs.\ $\mathrm{U}(d)$ distinction in finite matrix algebras.

\textbf{Source:} Section~8C, Table~XXVIII row 4.

%=============================================================================
\section{(c) Regular-Module Boost: $\varepsilon_{\mathrm{adj}} = d^2/(d^2-1) = 4096/4095$}
\label{sec:boost}
%=============================================================================

\subsection{The Forced Binary Fork}

The microsector construction requires choosing a finite Hilbert space
$\mathcal{H}_F$ on which the gauge action is defined. There are exactly
two canonical choices:

\paragraph{Branch A: Regular module (survives).}
\begin{equation}
\mathcal{H}_F := A = M_d(\mathbb{C}), \qquad \dim(\mathcal{H}_F) = d^2 = 4096.
\end{equation}
The gauge action is by inner derivations: $\mathrm{ad}_a(X) = [a, X]$.

\paragraph{Branch B: Fermion representation (falsified).}
\begin{equation}
\mathcal{H}_F := \mathbb{C}^d, \qquad \dim(\mathcal{H}_F) = d = 64.
\end{equation}
The gauge action is by fundamental matrix multiplication.

\subsection{Trace Conversion and the BOOST Factor}

\paragraph{UV (democratic) trace.}
The natural normalization for the spectral action trace on $\mathcal{H}_F = M_d(\mathbb{C})$
is the democratic normalization per matrix degree of freedom:
\begin{equation}
\mathrm{tr}_{\rm dem}(\cdot) := \frac{1}{d^2}\,\mathrm{Tr}_{\mathcal{H}_F}(\cdot).
\end{equation}

\paragraph{Physics (per-generator) trace.}
When reporting the gauge kinetic term in canonical $\mathfrak{su}(d)$ generator
normalization:
\begin{equation}
\mathrm{tr}_{\mathfrak{su}}(\cdot) = \frac{1}{d^2 - 1}\,\mathrm{Tr}_{\mathfrak{su}}(\cdot).
\end{equation}

\paragraph{The conversion factor is forced:}
\begin{equation}
\boxed{\varepsilon_{\mathrm{adj}}^{(A)} = \frac{d^2}{d^2 - 1} = \frac{4096}{4095}
= 1.000244\ldots \qquad \text{(Branch A: BOOST)}}
\end{equation}

This arises because $M_d(\mathbb{C})$ decomposes under the adjoint action as
$\mathfrak{su}(d) \oplus \mathbb{C}\cdot\mathbf{1}$, i.e., the $(d^2-1)$-dimensional
traceless part plus the 1-dimensional trace (identity). The democratic trace
averages over all $d^2$ directions, but only $d^2 - 1$ are physical gauge
directions. The ratio $d^2/(d^2-1)$ corrects for this overcounting.

\paragraph{Branch B gives the inverse:}
\begin{equation}
\varepsilon_{\mathrm{adj}}^{(B)} = \frac{d^2 - 1}{d^2} = \frac{4095}{4096}
= 0.999756\ldots \qquad \text{(Branch B: DROP)}
\end{equation}

\textbf{Source:} Section~8C, Eqs.~(8C.8)--(8C.10).

\subsection{Why Branch A Survives and Branch B Is Falsified}

Holding all other ingredients fixed:
\begin{center}
\renewcommand{\arraystretch}{1.3}
\begin{tabular}{lccc}
\toprule
Branch & Factor & $\alpha^{-1}$ & Residual (ppm) \\
\midrule
\textbf{A (regular-module)} & $4096/4095$ & $137.03599985$ & $-0.006$ \\
B (fermion-rep) & $4095/4096$ & $137.03014445$ & $+42.7$ \\
\midrule
Experimental & --- & $137.035999084$ & --- \\
\bottomrule
\end{tabular}
\end{center}

Branch B's 43~ppm deficit \textbf{cannot} be rescued by:
\begin{itemize}
\item Adjusting the U(1)/non-Abelian mixing weights ($w$): would require
  $\Delta w/w = -200\%$
\item Adjusting $g_F$: would require $\Delta g_F/g_F = +200\%$
\item Using $a_6$ correction: would require tuning cutoff moments to
  $f_{-2}/f_0 \sim 10^{-3}$, violating the no-knobs policy
\end{itemize}

Under the no-hidden-knobs policy, only Branch A survives.

\textbf{Source:} Section~8C, Sec.~\ref*{sec:alpha-decision-rule} and
Table~XXVI.

%=============================================================================
\section{(d) The Explicit Formula Connecting to $\kappa_{U(1)}$}
\label{sec:kappa}
%=============================================================================

\subsection{The Master Formula}

The electromagnetic coupling emerges from electroweak mixing. The DFD stiffness
functional for gauge sector $r$ gives:
\begin{equation}
g_1^2 = \frac{1}{\kappa_{U(1)}}, \qquad g_2^2 = \frac{4}{\kappa_{SU(2)}}.
\end{equation}

Electroweak mixing ($1/e^2 = 1/g_1^2 + 1/g_2^2$) combined with the Frame
Stiffness Theorem ($\kappa_{U(1)}/\kappa_{SU(2)} = 1/2$) yields:
\begin{equation}
\frac{1}{e^2} = \kappa_{U(1)} + \frac{2\kappa_{U(1)}}{4}
= \frac{3}{2}\,\kappa_{U(1)}.
\end{equation}

Therefore:
\begin{equation}
\boxed{\alpha^{-1} = \frac{4\pi}{e^2} = 6\pi\,\kappa_{U(1)}.}
\end{equation}

\textbf{Source:} Alpha\_Inverse\_Formula.tex, Sec.~2, Eq.~(5).

\subsection{How the Toeplitz Factors Enter $\kappa_{U(1)}$}

The stiffness $\kappa_{U(1)}$ is computed from the $a_4$ Seeley--DeWitt
coefficient of the spectral action on the Toeplitz-truncated internal geometry.
The computation incorporates all eight Table~XXVIII components:

\begin{enumerate}
\item \textbf{Connection Laplacian} $P = -g^{ij}\nabla_i\nabla_j$ with parallel
  K\"ahler curvature (no derivative ambiguities)

\item \textbf{Plateau cutoff} with $f_{-2} = f_{-4} = \cdots = 0$ (eliminates
  $a_6$ and higher as hidden tuning knobs)

\item \textbf{Toeplitz truncation} at level $m = k+3$, giving $d = k+4 = 64$
  and spectral cutoff $\Lambda^3 = k(d-1)/d = 60 \times 63/64 = 59.0625$

\item \textbf{SM content:} hypercharge weighting
  $w = N_{\mathrm{species}}/(g_F \cdot \mathrm{Tr}(Y^2)) = 7/(8 \times 10)
  = 7/80$

\item \textbf{Regular-module BOOST:}
  $\varepsilon_{\mathrm{adj}} = d^2/(d^2-1) = 4096/4095$
\end{enumerate}

The result:
\begin{equation}
\kappa_{U(1)}^{\mathrm{pred}} = 7.26999\ldots
\end{equation}

Therefore:
\begin{equation}
\alpha^{-1} = 6\pi \times 7.26999\ldots = 137.03599985.
\end{equation}

\textbf{Source:} Alpha\_Inverse\_Formula.tex, Sec.~5.3.

\subsection{The Approximate Closed-Form Formula}

While no exact single algebraic expression exists (the full computation
requires numerical evaluation of the $a_4$ coefficient), an approximate
closed form captures the structure to 85~ppm:
\begin{equation}
\alpha^{-1} \approx \frac{35\pi}{3}\;\beta_{U(1)}\;\frac{d-1}{d}\;
\frac{d^2}{d^2-1},
\end{equation}
where $\beta_{U(1)} = 3.7969$ is the Chern--Simons weighted average at
$k_{\max} = 60$. This factors as:
\begin{equation}
\frac{35\pi}{3} = 4\pi \times \underbrace{\frac{5}{2}}_{\text{EW mixing}}
\times \underbrace{\frac{7}{6}}_{\text{spectral triple}},
\end{equation}
where $5/2 = 1 + R_W/4$ with Wilson ratio $R_W = 6$, and
$7/6 = N_{\mathrm{species}}/(N_{\mathrm{gen}} \cdot n_2) = 7/(3 \times 2)$.

Numerically:
\begin{equation}
\frac{35\pi}{3} \times 3.7969 \times \frac{63}{64} \times \frac{4096}{4095}
= 137.024 \qquad \text{(85~ppm residual)}.
\end{equation}

\textbf{Source:} Alpha\_Inverse\_Formula.tex, Sec.~4.

%=============================================================================
\section{The Toeplitz Truncation in Other DFD Results}
\label{sec:other}
%=============================================================================

The same $d = k+4$ Toeplitz truncation machinery appears in several other
DFD derivations:

\subsection{The $\alpha^{57}$ Cosmological Constant Exponent}

The primed-determinant scaling on the finite Toeplitz state space
$\mathcal{H}$ of dimension $k_{\max} = 60$ with $\dim\ker(\mathcal{K})
= N_{\mathrm{gen}} = 3$ gives:
\begin{equation}
\varepsilon(\alpha^{-1}) = \alpha^{k_{\max} - N_{\mathrm{gen}}} = \alpha^{57}.
\end{equation}
This is identified with $G\hbar H_0^2/c^5$ via the observer dictionary.

\textbf{Source:} Appendix~O.

\subsection{The Majorana Scale $M_R = M_P\alpha^3$}

The Toeplitz-quantized singlet sector has $\dim(\mathcal{H}_{\nu_R})
= N_{\mathrm{gen}} = 3$. Determinant scaling gives:
\begin{equation}
M_R = M_P \times \alpha^{N_{\mathrm{gen}}} = M_P \alpha^3
= 4.74 \times 10^{12}\,\mathrm{GeV}.
\end{equation}

\textbf{Source:} Appendix~P, Theorem~P.3.

\subsection{Neutrino Mass Ratios}

The $\mathbb{C}P^1$ Toeplitz truncation with $m_\nu = a + 4$ (where
$a = 9$ from minimal-padding) gives finite dimension
$d_\nu = m_\nu + 1 = a + 5 = 14$. Channel fractions like $2/13$ and
$3/11$ from this truncation fix the neutrino mass ratios as powers of
$\alpha$.

\textbf{Source:} Appendix~X.

%=============================================================================
\section{Summary: The Four Mechanisms}
\label{sec:summary}
%=============================================================================

\begin{tcolorbox}[colback=blue!5!white, colframe=blue!50!black,
  title=Toeplitz Truncation at $d=64$: Complete Mathematical Anatomy]

\textbf{(a) Spectrum modification:}
The continuous $\mathbb{C}P^2 \times S^3$ is replaced by a $64$-dimensional
matrix algebra $M_{64}(\mathbb{C})$. The spectral action is computed via
the $a_4$ Seeley--DeWitt coefficient of the connection Laplacian on this
finite geometry, with plateau cutoff eliminating higher heat-kernel terms.
The spectral cutoff is $\Lambda^3 = 59.0625$.

\textbf{(b) Traceless projection $(d-1)/d = 63/64$:}
Of $d = 64$ levels, only $d-1 = 63$ contribute to traceless $\mathrm{SU}(d)$
gauge fields. The determinant channel is removed. This enters the spectral
cutoff: $\Lambda^3 = k \times (d-1)/d = 60 \times 63/64$.

\textbf{(c) Regular-module boost $d^2/(d^2-1) = 4096/4095$:}
The finite Hilbert space is $\mathcal{H}_F = M_d(\mathbb{C})$ (dimension
$d^2 = 4096$), not $\mathbb{C}^d$ (dimension $d = 64$). Converting from
democratic UV trace ($1/d^2$ per direction) to physics trace ($1/(d^2-1)$
per $\mathfrak{su}(d)$ generator) forces $\varepsilon_{\mathrm{adj}}
= 4096/4095$. This is the \textbf{forced binary fork}; Branch B
($4095/4096$) is falsified at 43~ppm.

\textbf{(d) Connection to $\kappa_{U(1)}$:}
The master formula is $\alpha^{-1} = 6\pi\,\kappa_{U(1)}$, where
$\kappa_{U(1)} = 7.26999\ldots$ is computed numerically from the $a_4$
coefficient incorporating all eight locked inputs (Table~XXVIII).
No single closed-form algebraic expression exists at sub-ppm precision;
the approximate formula
$\alpha^{-1} \approx (35\pi/3)\,\beta\,(63/64)\,(4096/4095) = 137.024$
captures the structure to 85~ppm.

\textbf{Final result:}
\begin{equation}
\alpha^{-1} = 6\pi \times 7.26999\ldots = 137.03599985
\qquad\text{(residual: $-0.006$~ppm from experiment)}.
\end{equation}
\end{tcolorbox}

%=============================================================================
\section*{Source Files}
%=============================================================================

\begin{itemize}
\item \texttt{section\_alpha\_locked.tex} --- Section~8C: Convention-locked
  $\alpha$ from the microsector (primary source for all four mechanisms)
\item \texttt{appendix\_K.tex} --- Appendix~K: Microsector physics, Bridge
  Lemma, $k_{\max} = 60$ derivation
\item \texttt{appendix\_O\_alpha57.tex} --- Appendix~O: Primed-determinant
  scaling and $\alpha^{57}$ exponent
\item \texttt{appendix\_P\_kalpha\_MR.tex} --- Appendix~P: Toeplitz scaling
  for clock coupling and Majorana scale
\item \texttt{appendix\_X\_neutrino.tex} --- Appendix~X: Toeplitz truncation
  in neutrino mass ratios
\item \texttt{Alpha\_Inverse\_Formula.tex} --- Cross-reference document:
  explicit derivation chain and approximate closed form
\end{itemize}

\end{document}
